\section{An output-size obstruction for compressed cuts}
\label{sec:compressed-cut-output}

An efficient compressed cutting theorem must be translated into an
explicit output contract.  Merely storing normal coordinates in binary
does not make an expanded component list efficient.

\begin{proposition}[Parallel-disc output lower bound]
\label{prop:parallel-disc-output}
There is a fixed one-tetrahedron triangulated 3-ball and a family of
properly embedded orientable normal surfaces $S_W$, for integers $W\ge1$,
with normal-coordinate input length $O(\log(W+1))$, such that cutting
along $S_W$ produces $W+1$ connected components.  Between successive
normal discs are $W-1$ pairwise disjoint product slabs.  Consequently an
algorithm that explicitly lists all cut components, or all these
individual product slabs, requires $\Omega(W)$ output operations.
\end{proposition}

\begin{proof}
Choose a vertex of one tetrahedron and take $W$ pairwise disjoint,
nested, normally isotopic triangles linking that vertex.  Their normal
vector has one coordinate $W$ and six zero coordinates.  Each triangle
is a properly embedded disc.  In a collar of one such disc, the
successive triangles divide the collar into $W-1$ copies of
$D^2\times I$.  The two remaining regions, one towards the vertex and
one away from it, are also balls.  Therefore the cut manifold has
$W+1$ components.  Writing one record for each component or slab has
the asserted lower bound.  Taking $W=2^k$ makes this exponential in
the binary coordinate length.
\end{proof}

This does not contradict a compressed-cutting theorem.  It shows that
its concrete interface must support aggregated multiplicities,
symbolic families, or a stated connectedness restriction.  A suitable
record for this example is one ordinary product type with binary
multiplicity $W-1$, together with the two exceptional end regions.
It is insufficient to advertise polynomial time in $\log W$ while
returning a Python list with one element per product region.

More generally, if $S$ is a connected two-sided normal surface with
normal vector $v$, then $Wv$ is represented by $W$ parallel copies of
$S$.  With compatible collar markings,
\[
 M\mathbin{\backslash\!\backslash}(WS)
 \cong
 \bigl(M\mathbin{\backslash\!\backslash}S\bigr)
 \sqcup
 \coprod_{i=1}^{W-1}(S\times I).
\]
Thus multiplicity extraction can be an exact preprocessing operation
when connectedness and two-sidedness of the base surface are certified.
Dividing coordinates by their greatest common divisor alone does not
certify either condition.  For a one-sided connected surface, doubling
its normal vector describes the boundary of a regular neighborhood,
which can be connected; blindly interpreting that factor of two as
two disjoint copies would be incorrect.

In the July 2026 hierarchy paper, Theorem 14.4 (p.~51) gives the
compressed cut information, whereas Proposition 14.2 (p.~50) adds
essential-pattern and irreducibility hypotheses for the subsequent
simplification.  These are different contracts.  In particular,
an undecided unknot exterior does not supply an essential initial
pattern as an available assumption.  The representation lower bound
above is an implementation obligation, not an assertion that the
literature's compressed theorem is false.
